

%%%%%%%%%%%%%%%%%%%%%%%%%%%%%%%%%%%%%%%%%%%%%%%%%%%%%%%%%%%%%%%%%%%%
%全文需要的包
%不包括只需要在开头进行一次调用的，后面都再也用不到的包，
%那些都遵循靠近使用处的原则，在format内声明


%\usepackage[UTF8]{ctex} %中文支持

\usepackage{braket}
\usepackage{import}

\usepackage{color} % 支持彩色

\usepackage{multirow}
\usepackage{tikz}
\usepackage{pgf}
\usepackage{diagbox}
\usepackage{amsmath} % AMSLaTeX宏包 用来排出更加漂亮的公式
\usepackage{amsthm}
\usepackage{amssymb} % 数学符号生成命令
\usepackage{siunitx} %单位、数字宏包 常用的有 \num，\SI{10}{\hertz}，\SIrange{10}{100}{\hertz}
%\usepackage[colorlinks=false, pdfborder={0 0 0}]{hyperref} %hyperref 非常强大，你可以有非常多的可能性，其中最突出的特色是超链接。当引用一幅图的时候，引用与图形形成了链接，当你点击引用的地方，它会跳转到链接的图片处。并且 hyperref 可以让你插入 PDF 元数据到你的最终文档中。注意：作为一个经验法则，你应该在导言区的最后加入这个宏包，在所有宏包之后。也存在少数例外的情况：比如，本文提到的 cleveref 宏包，cleveref宏包应该在 hyperref 之后。更多的例外情况可以参看：pkg load after hyperref。
%\usepackage[amsmath,thmmarks,hyperref]{ntheorem}  % 定理类环境宏包，其中 amsmath 选项用来兼容 AMS LaTeX 的宏包


\usepackage{graphicx}
\graphicspath{{figures/}}

%\usepackage{caption,titletoc,titlesec}%标题
%\usepackage{titletoc,titlesec}%标题

% 超链接
%\usepackage{cleveref} %这个宏包引入了 \cref 命令，当使用这个命令用于交叉引用的时候（而不是 \ref 或者 \eqref），根据引用的不同，它会自动添加一个单词前缀，引用 figure 环境，它会自动添加 “fig.”,而对于 equation 环境，它会自动添加 “eq.”。

